\section{Weighted information gain}

In the standard Gaussian process setting, the maximum information gain after $T$ observations is
\begin{equation*}
\Gamma_T
= \frac{1}{2}
\max_{x_1,\dots,x_T}
\sum_{t=1}^T
\log\!\left(
1+\frac{\sigma_{t-1}^2(x_t)}{\sigma^2}
\right),
\end{equation*}
where $\sigma_{t-1}^2(x_t)$ denotes the posterior variance before observing $x_t$.

Motivated by the preceding discussion, we introduce the weighted sequential objective
\begin{equation*}
\Gamma_T^w
= \frac{1}{2}
\sum_{t=1}^T
w(x_t)^2
\log\!\left(
1+\frac{\sigma_{t-1}^2(x_t)}{\sigma^2}
\right).
\end{equation*}

The weights act externally on the marginal gains while leaving the Gaussian process posterior unchanged. The motivation is that if $w$ is concentrated on a sufficiently small subset of the domain, then $\Gamma_T^w$ may admit sharper bounds than the ordinary information gain.
For example, if
\begin{equation*}
w(x)=c\,\mathbf{1}_A(x),
\end{equation*}
then
\begin{equation*}
\Gamma_T^w
= \frac{c^2}{2}
\sum_{t:\,x_t\in A}
\log\!\left(
1+\frac{\sigma_{t-1}^2(x_t)}{\sigma^2}
\right),
\end{equation*}
which resembles a transductive information gain restricted to $A$.

\subsection{Loss of submodularity for weighted objective}

For ordinary information gain, the set function
\begin{equation*}
F(S)=I(f_S;y_S)
\end{equation*}
has marginal gains
\begin{equation*}
F(S\cup\{x\})-F(S)
=
\frac12
\log\!\left(
1+\frac{\sigma_S^2(x)}{\sigma^2}
\right),
\end{equation*}
which satisfy diminishing returns. This submodularity is what allows standard greedy approximation arguments.

The weighted objective does not inherit this structure. Writing formally
\begin{equation*}
F_w(x_1,\dots,x_T)
= \frac{1}{2}
\sum_{t=1}^T
w(x_t)^2
\log\!\left(
1+\frac{\sigma_{t-1}^2(x_t)}{\sigma^2}
\right),
\end{equation*}
the contribution of a point depends both on its weight and on its position in the sequence through the posterior variance. As a result, the objective is order-dependent and thus cannot be expressed as a submodular set function.This prevents the direct application of the standard TAL framework.